\documentclass[aspectratio=169,11pt]{beamer}

% ---- clean, modern, dependency-light look ----
\usetheme{default}
\usefonttheme{professionalfonts}
\usepackage[T1]{fontenc}
\usepackage{lmodern}
\usepackage[utf8]{inputenc}
\usepackage{amsmath,amssymb}
\usepackage{booktabs}
\usepackage{graphicx}
\usepackage{xcolor}
\graphicspath{{../results/plots/}}

\setbeamertemplate{navigation symbols}{}

% palette
\definecolor{ink}{RGB}{28,38,56}
\definecolor{accent}{RGB}{0,104,168}
\definecolor{accent2}{RGB}{192,57,43}
\definecolor{good}{RGB}{34,110,55}
\definecolor{soft}{RGB}{244,247,250}

\setbeamercolor{normal text}{fg=ink}
\setbeamercolor{frametitle}{fg=accent,bg=}
\setbeamercolor{title}{fg=ink}
\setbeamercolor{structure}{fg=accent}
\setbeamercolor{block title}{fg=white,bg=accent}
\setbeamercolor{block body}{fg=ink,bg=soft}
\setbeamerfont{frametitle}{series=\bfseries,size=\large}
\setbeamerfont{title}{series=\bfseries,size=\Large}
\setbeamertemplate{itemize item}{\color{accent}\textbullet}
\setbeamertemplate{itemize subitem}{\color{accent}--}

% flat frametitle with a thin rule, generous padding
\setbeamertemplate{frametitle}{%
  \vspace{6pt}\hspace{6pt}{\usebeamerfont{frametitle}\insertframetitle}%
  \vspace{2pt}\par%
  {\color{accent}\hrule height 1.5pt}%
  \vspace{-2pt}%
}
% minimal footer: short title (left) -- page (right)
\setbeamertemplate{footline}{%
  \hbox{}\vspace{-4pt}%
  \hspace{8pt}{\color{ink!55}\scriptsize CVaR-DQN Exam Scheduler}%
  \hfill{\color{ink!55}\scriptsize \insertframenumber/\inserttotalframenumber}%
  \hspace{8pt}\vspace{4pt}%
}

\newcommand{\hl}[1]{{\color{accent}\textbf{#1}}}
\newcommand{\win}[1]{{\color{good}\textbf{#1}}}
\newcommand{\bad}[1]{{\color{accent2}\textbf{#1}}}

\title{Risk-Sensitive RL for Study Scheduling}
\subtitle{A CVaR-DQN Approach}
\author{Kobiljon Muhammadov}
\institute{RL and Decision Making Under Uncertainty --- Spring 2026}
\date{}

\begin{document}

% ---------- 1. Title (~10s) ----------
{\setbeamertemplate{footline}{}
\begin{frame}
  \vfill
  {\color{accent}\rule{\textwidth}{2pt}}\\[4pt]
  {\usebeamerfont{title}\color{ink}\inserttitle}\\[2pt]
  {\large\color{accent}\insertsubtitle}\\[10pt]
  {\insertauthor \hfill \insertinstitute}\\[4pt]
  {\color{accent}\rule{\textwidth}{2pt}}
  \vfill
  \begin{center}\itshape\color{ink!70}
  Goal: schedule study days so the student does not fail \emph{any} exam.
  \end{center}
  \vfill
\end{frame}}

% ---------- 2. Problem (~50s) ----------
\begin{frame}{The Problem}
  A student has \hl{$D$ days} before exams. Each day: study \emph{one} subject, or rest.
  \begin{itemize}
    \item \hl{Forgetting:} unstudied subjects decay (Ebbinghaus curve).
    \item \hl{Energy:} random day-to-day; low energy $\Rightarrow$ learn less.
    \item \hl{Weights:} exams carry different credit $w_i$.
  \end{itemize}
  \vspace{0.8em}
  \begin{block}{The tension --- our research question}
    Maximising the \emph{average} score makes the optimal policy
    \bad{sacrifice a low-weight subject}. But a student cares about
    \win{not failing \emph{any} exam}. We need \hl{risk-sensitive} RL.
  \end{block}
\end{frame}

% ---------- 3. Environment (~50s) ----------
\begin{frame}{The Environment $=$ a Markov Decision Process}
  \hl{State} $s=(\mathbf{k},e,d)$: knowledge per subject, energy, day.\quad
  \hl{Action} $a$: study a subject, or rest.
  \vspace{0.6em}
  \[
  \underbrace{k_i \leftarrow \lfloor k_i e^{-\lambda_i}\rfloor}_{\textbf{forget}}\qquad
  \underbrace{k_a \leftarrow k_a + \lfloor \tfrac{e}{E}\,g_{\max}\rfloor}_{\textbf{learn}}\qquad
  \underbrace{r = w_i\,\sigma(k_i-\tfrac{K}{2})}_{\textbf{exam reward}}
  \]
  \vspace{0.4em}
  \begin{itemize}
    \item Sigmoid reward: knowing $K/2$ $\Rightarrow$ exactly half credit $0.5\,w_i$.
    \item Energy noise $\xi\!\sim\!\mathrm{Uniform}\{-1,0,+1\}$ is the \emph{only} randomness.
    \item Three scales: Small / Medium / Large ($\sim$7.5k / 328k / 11M states).
  \end{itemize}
\end{frame}

% ---------- 4. Core idea + figure (~65s) ----------
\begin{frame}{Core Idea: Optimise the Worst Case, Not the Average}
  \begin{columns}[T]
    \begin{column}{0.48\textwidth}
      \footnotesize
      \hl{DQN} learns one number: the mean.\\[0.3em]
      \hl{CVaR-DQN} learns the \emph{whole return distribution}
      ($M{=}32$ quantiles) and acts on the bad tail:
      \[
        a^{*}=\arg\max_a\ \
        \underbrace{\frac{1}{\lceil\alpha M\rceil}
        \sum_{i=1}^{\lceil\alpha M\rceil}\theta_{(i)}}
        _{\textstyle \mathrm{CVaR}_\alpha}
      \]
      $\theta_{(i)}$: $i$-th \emph{worst} return\,;\;
      $\lceil\alpha M\rceil$: \# tail quantiles ($3$ of $32$)\,;\;
      $\alpha{=}0.1$.\\[0.4em]
      $\Rightarrow$ \textbf{average of the worst 10\% of outcomes.}\\[0.7em]
      DQN: ``best on \emph{average}?''\\
      CVaR-DQN: ``best when I'm \emph{unlucky}?''
    \end{column}
    \begin{column}{0.52\textwidth}
      \vspace{1.5em}
      \includegraphics[width=\textwidth]{cvar_concept.png}\\[0.4em]
      \centering\footnotesize
      Failing one exam \emph{is} a worst-case event.
    \end{column}
  \end{columns}
\end{frame}

% ---------- 5. Methods (~40s) ----------
\begin{frame}{Methods Compared}
  \begin{columns}[T]
    \begin{column}{0.5\textwidth}
      \textbf{Baselines}
      \begin{itemize}
        \item 3 heuristics (Uniform, MUF, LKF)
        \item Value Iteration --- exact optimum (Small)
        \item Tabular Q-learning
        \item DQN $+$ Double DQN
      \end{itemize}
    \end{column}
    \begin{column}{0.5\textwidth}
      \textbf{Contributions}
      \begin{itemize}
        \item \hl{CVaR-DQN} --- risk-sensitive
        \item \hl{SC-CVaR-DQN} --- per-subject constrained \emph{(pilot)}
      \end{itemize}
      \vspace{1.2em}
      {\footnotesize 5 seeds $\times$ 5k train / 1k eval episodes each.}
    \end{column}
  \end{columns}
\end{frame}

% ---------- 6. Main result (~60s) ----------
\begin{frame}{Main Result: the Mean--Risk Trade-off}
  \begin{center}
  \begin{tabular}{llrr}
    \toprule
    Config & Agent & Mean Return & CVaR$_{10\%}$ \\
    \midrule
    Small  & Value Iteration & \textit{5.00} & \textit{3.44} \\
    Small  & DQN             & 4.97 & 3.37 \\
    \midrule
    Medium & DQN             & \textbf{4.86} & 2.86 \\
    Medium & CVaR-DQN        & 4.50 & \bad{3.19} \\
    \midrule
    Large  & DQN             & \textbf{11.47} & 8.25 \\
    Large  & CVaR-DQN        & 10.62 & \bad{8.27} \\
    \bottomrule
  \end{tabular}
  \end{center}
  \vspace{0.4em}
  \begin{block}{The contribution, in one line}
    DQN $4.97\!\approx\!$ VI optimum $5.00$ (it \emph{found} the optimum).
    On Medium: trade \bad{7\% average} for \win{12\% better worst-case}
    --- the finance mean--risk principle, realised.
  \end{block}
\end{frame}

% ---------- 7. Pilot + figure (~65s) ----------
\begin{frame}{Pilot: Subject-Constrained CVaR-DQN (honest findings)}
  Constrain \emph{each subject $i$} ($Z_i$ = its return, $T_i$ = its pass bar):
  \vspace{-0.2em}
  \[
    a^{*}=\arg\max_a\ \
    \underbrace{\sum_i\mathbb{E}[Z_i]}_{\text{maximise total reward}}
    \ -\ \lambda
    \underbrace{\sum_i\big[\max(0,\ T_i-\mathrm{CVaR}_\alpha(Z_i))\big]^2}
    _{\text{penalty if any subject's worst case is below its bar}}
  \]
  \begin{columns}[T]
    \begin{column}{0.50\textwidth}
      \footnotesize
      \vspace{0.3em}
      \win{Win:} Medium Physics failure
      $30\%\!\to\!\mathbf{8\%}$ (4--6$\times$).\\[0.5em]
      \bad{Honest losses:} lower aggregate CVaR; worse under shift;
      cannot help \emph{infeasible} subjects.\\[0.7em]
      \textbf{Key insight:} the lowest-weight subject fails 100\%
      \emph{even for Value Iteration} --- an \hl{environment property},
      not an algorithm flaw.
    \end{column}
    \begin{column}{0.50\textwidth}
      \includegraphics[width=\textwidth]{per_subject_failure.png}
    \end{column}
  \end{columns}
\end{frame}

% ---------- 8. Takeaways (~30s) ----------
\begin{frame}{Takeaways}
  \begin{enumerate}
    \item \hl{CVaR-DQN works:} consistently better worst-case than DQN,
          at a small cost in mean --- the mean--risk trade-off.
    \item \hl{Energy noise is everything:} the constant-energy ablation
          makes DQN $=$ CVaR-DQN. It is the sole source of risk.
    \item \hl{The pilot is an honest mixed result:} one clear win, clear
          failure modes, one concrete next step --- \emph{adaptive thresholds}.
  \end{enumerate}
  \vfill
  \begin{center}
    {\Large\color{accent}\textbf{Thank you --- Questions?}}\\[0.4em]
    {\footnotesize\texttt{github.com/muhammadov-q/cvar-dqn-exam-scheduler}}
  \end{center}
\end{frame}

% ===================== BACKUP =====================
\appendix
{\setbeamertemplate{footline}{\hfill{\color{ink!40}\scriptsize backup}\hspace{8pt}\vspace{4pt}}

\begin{frame}{Backup: Ablation --- proof the mechanism is \emph{risk}}
  \begin{center}\small
  \begin{tabular}{lrrrr}
    \toprule
    Ablation (Medium) & DQN mean & CVaR mean & DQN CVaR & CVaR-DQN CVaR\\
    \midrule
    Baseline         & 4.79 & 4.50 & 2.91 & \bad{3.19}\\
    \textbf{Constant energy} & \textbf{7.01} & \textbf{7.01} & \textbf{7.01} & \textbf{7.01}\\
    No forgetting    & 5.63 & 5.59 & 4.34 & \bad{4.95}\\
    $\lambda{=}0.4$  & 4.57 & 4.36 & \textbf{3.03} & 2.86\\
    \bottomrule
  \end{tabular}
  \end{center}
  \vspace{0.5em}
  Remove the energy noise $\Rightarrow$ the environment is deterministic
  $\Rightarrow$ there is \emph{no risk} $\Rightarrow$ DQN and CVaR-DQN become
  \hl{identical} (both 7.01). This \emph{proves} the difference between them
  is genuinely risk-sensitivity, not training noise.
\end{frame}

\begin{frame}{Backup: Why $M=32$ quantiles?}
  \begin{itemize}
    \item Standard distributional-RL default; sets distribution resolution.
    \item With $\alpha{=}0.1$: $\lceil 0.1{\times}32\rceil = 3$ tail quantiles
          --- enough for a \emph{stable} CVaR average, not a single noisy point.
    \item Too small ($M{=}10\Rightarrow$ 1 tail point $=$ noisy, that's VaR not CVaR);
          too big ($M{=}200$) $=$ overkill, slower, bigger per-subject net.
    \item Not exhaustively swept --- a principled default, not a tuned optimum.
  \end{itemize}
\end{frame}

\begin{frame}{Backup: Why the \emph{quadratic} Lagrangian penalty?}
  \begin{itemize}
    \item Hard constraint $\Rightarrow$ feasible set empty early in training
          $\Rightarrow$ discontinuous Bellman target $\Rightarrow$ unstable.
    \item Lagrangian relaxation makes the objective continuous.
    \item \emph{Quadratic}, not linear: when all actions violate, a linear
          penalty is a near-constant offset that doesn't differentiate
          actions; squaring amplifies relative differences so the
          least-bad action wins.
  \end{itemize}
\end{frame}

\begin{frame}{Backup: CVaR vs Variance vs VaR}
  \begin{itemize}
    \item \textbf{Variance} punishes good surprises too; CVaR only the downside.
    \item \textbf{VaR} $=$ just the cutoff quantile; ignores how bad the tail
          is beyond it; non-convex, hard to optimise.
    \item \textbf{CVaR} $=$ average of the whole tail; coherent, convex,
          matches ``how bad on average when it goes wrong''.
  \end{itemize}
\end{frame}
}

\end{document}
